\documentclass[12pt]{article}
\usepackage[utf8]{inputenc}
\usepackage[T1]{fontenc}
\usepackage{hyperref}
\title{Xipaya Language Resources for K-12 Curriculum Development}
\author{Compiled by Claude Sonnet 5.5 (claude-sonnet-5-5) for Andr\'as Kornai}
\date{}
\begin{document}
\maketitle

\section{Introduction}
Xipaya (also Shipaya, Xipaia, Chipaya; ISO 639-3 xiy; Glottocode xipa1240) is an endangered Tupian language of the Juruna branch, spoken in Par\'a, Brazil, and closely related to Juruna. Wikipedia reports one remaining native speaker (as of 2024), an elderly woman in Altamira. An older mirror gives two. Most Xipaya people speak Portuguese now. Some older people still know the language but do not use it daily. The search found no school programme or teaching materials.

\section{Primary Educational Resources}
The search found no primers, textbooks or school programmes for Xipaya.

\section{Linguistic and Reference Materials}
\begin{itemize}
\item Wikipedia, ``Xipaya language'': \url{https://en.wikipedia.org/wiki/Xipaya_language}
\item Povos Ind\'igenas no Brasil (Instituto Socioambiental), Xipaya: \url{https://pib.socioambiental.org/en/povo/Xipaya}
\item Paper in \emph{Moara} (UFPA) turned up in the search; I did not confirm its subject beyond the results: \url{https://periodicos.ufpa.br/index.php/moara/article/download/10906/7505}
\item Cristina Martins Fargetti and colleagues have published on Juruna and Xipaya, including a 2021 comparison of vowel systems: \url{https://unep.academia.edu/CristinaMartinsFargetti}
\end{itemize}

\section{Religious and Community Resources}
The search found none, other than the community profile at Povos Ind\'igenas no Brasil above.

\section{Digital and Audio Resources}
\begin{itemize}
\item ELDP collection ``Documentation of Urgently Endangered Tupian Languages'', covering five endangered Brazilian languages: \url{https://eldp.access.preservica.com/dk0233/}. I did not confirm that Xipaya is among the recordings, though the search associated the collection with it.
\item Endangered Languages Project entry: \url{https://endangeredlanguages.com/lang/1001?hl=pt-br}
\end{itemize}

\section{Curriculum Development Context}
With at most one or two elderly speakers, any programme would be documentation-led reclamation. It would draw on the ELDP archive and academic descriptions, and it would be done with the Xipaya community in Altamira. Such a programme is at an early stage and does not yet have school materials.
\end{document}